\chapter{East Chadic (16/36)}

Information is available on the grammars of 16 of 36 East Chadic languages. This branch is divided into two subbranches. The East Chadic A languages are spoken in the south of Chad, mainly in the Chari-Logone area. The East Chadic B languages are spoken mainly in the Guera region of Chad. 

\section{East Chadic A (Chari-Logone) (5/14)}

Information on the grammars of East Chadic A languages is limited to five languages, of which two report having any directional extensions. 

%As seen in Table \ref{tab:EastA}, only two East Chadic A languages have been shown to have any directional extensions. Both Lele and Kera have a ventive extensions. The two forms do not immediately appear to be cognate. 

%\begin{table}[H]
%	\caption{East Chadic A directional morphology}
%	\label{tab:EastA}
%	\begin{tabular}{lllllll} % add l for every additional column or remove as necessary
%		\lsptoprule
%		language & VENT & ITIVE  & UP & DOWN & INTO & OUT  \\ %table header
%		\midrule
%		Ndam (ndam1251) &  &  & & & &  \\	
%		Somrai (somr1248) &  & 	& & & &  \\	
%		Lele (lele1276) & \textit{jè} & & & & &  \\	
%		Kera (kera1255) & \textit{-dà} & & & & &  \\	
%		Kwang (kwan1285) &  & 	& & & &  \\	
%		\lspbottomrule
%	\end{tabular}
%\end{table}

%\begin{table}[H]
%	\caption{East Chadic A directional verbs}
%	\label{tab:EastAV}
%	\begin{tabular}{lllllll} % add l for every additional column or remove as necessary
%		\lsptoprule
%		language & VENT & ITIVE  & UP & DOWN & INTO & OUT  \\ %table header
%		\midrule
%		Ndam (ndam1251) & \textit{ugʌ᷆}  &  & \textit{əga᷆} & \textit{jə̀bʌ̀} & \textit{ɗúgéy} &  \textit{ɗə́gâ} \\	
%		Somrai (somr1248) &  \textit{háde}   & \textit{há}	& \textit{ájìl} & \textit{sə́} & \textit{də̄mə̄} & \textit{ganda}  \\	
%		Lele (lele1276) & \textit{è} & \textit{èjè} &                  & \textit{jìrmí} & \textit{do'b} & \textit{an}  \\	
%		Kera (kera1255) & \textit{bi} & \textit{ɗé} & \textit{lí} & \textit{lí wə́ra} & \textit{kélé} & \textit{kél wə́ra} \\	
%		Kwang (kwan1285) & \textit{ɗéːna}  & \textit{ɗe} & \textit{ale} & \textit{jìrpe, jàːpe} & \textit{gálkájaŋ}, \textit{báyé} & \textit{gale}, \textit{gáljóŋo} \\	
%		\lspbottomrule
%	\end{tabular}
%\end{table}

%Details about the expression of directed CAM is only available for two East Chadic A languages. Lele uses a prepositional phrase with a motion verb, while Kera uses a verb glossed `take' with a ventive suffix. Most East Chadic A languages have a similar verb root meaning `trade' or `buy' and, in the three cases where data is available, each language uses a different morph for marking a SELL meaning using the same verb root, as seen in Table \ref{tab:EastABUY}.

%\begin{table}[H]
%	\caption{BUY and SELL in East Chadic A}
%	\label{tab:EastABUY}
%	\begin{tabular}{lll} % add l for every additional column or remove as necessary
%		\lsptoprule
%		language & TRADE/BUY  & SELL   \\ %table header
%		\midrule
%		Ndam (ndam1251) & \textit{kə́lə̂}  &  \\	
%		Somrai (somr1248) &   \textit{kələ̂} &    \\	
%		Lele (lele1276) & \textit{kil}    &  \textit{kil ... cáaní} \\	
%		Kera (kera1255) &  \textit{gùsí} & \textit{gùsí wə̀ra} \\	
%		Kwang (kwan1285) &  \textit{kèle} & \textit{kèlè sə́ŋ} \\	
%		\lspbottomrule
%	\end{tabular}
%\end{table}

\subsection{East Chadic A1 (2/6)}

Relevant information is available for two of the six East Chadic A1 group languages: Ndam and Somrai. No relevant publications were found describing the grammar of Gadang (gada1262), Mire (mire1238) or Sarua (saru1245) and only a very limited description (with no analysis of verbal morphology) is available for Tumak (tuma1260).

\paragraph*{Ndam (ndam1251)} \label{ndam}

Information on the grammar of Ndam is from a thesis on the phonology and morphology of the language \citep{brossMaterialenZurSprache1988}.

\subparagraph{Directional verbs}~

\begin{table}[H]
	\caption{Ndam directional verbs}
	\label{tab:ndamV}
	\begin{tabular}{lll} % add l for every additional column or remove as necessary
		\lsptoprule
		DIR & verb & gloss   \\ %table header
		\midrule
		VENT &\textit{ugʌ̀} & `\textit{venir} [come]' \\
		INTO& \textit{ɗúgéy} & `\textit{entrer} [enter]' \\
		OUT & \textit{ɗə́gâ} & `\textit{sortir} [go out]' \\
		UP & \textit{əga} & `\textit{gravir}, \textit{monter} [climb, go up]' \\
		DOWN & \textit{jə̀bʌ̀} & `\textit{descendre} [descend]' \\
		\lspbottomrule
	\end{tabular}
\end{table}

\subparagraph{Directional extensions}

\citet{brossMaterialenZurSprache1988} does not describe any motion-related verbal morphology.

\subparagraph{Caused accompanied motion (CAM)}

Caused accompanied motion may be expressed by the verb \textit{ugéy} `\textit{amener} (\textit{apporter}) [bring (carry)]' or \textit{́aay} `\textit{ramener} (\textit{apporter}) [bring back (carry)]' but it is not clear how directed CAM is expressed \citep[66, 90]{brossMaterialenZurSprache1988}. 

\subparagraph{Buy-sell verbs}

\citet[81]{brossMaterialenZurSprache1988} lists the verb \textit{kə́lə̂} `\textit{vendre}, \textit{acheter} [sell, buy]' but does not give any sense of how the notions `buy' and `sell' could be distinguished. 

\paragraph*{Somrai (somr1248)} \label{somr}

Information on the grammar of Somrai is from a publication of texts in the language \citep{jungraithmayrPresentationDunConte1978}, a sketch of the verbal system published as a chapter in an edited volume \citep{barreteauVerbeSibine1982} and a thesis on information structure \citep{millerInformationStructureSoumraye2020}.

\subparagraph{Directional verbs}~

\begin{table}[H]
	\caption{Somrai directional verbs}
	\label{tab:somrV}
	\begin{tabular}{lll} % add l for every additional column or remove as necessary
		\lsptoprule
		DIR & verb & gloss   \\ %table header
		\midrule
		ITIVE & \textit{há} & `\textit{aller} [go]' \\
		VENT & \textit{háde} & `\textit{aller vers ici}, \textit{venir} [come]' \\
		INTO& \textit{də̄mə̄} & `\textit{entrer} [enter]' \\
		OUT & \textit{ganda} & `come out' \\
		UP & \textit{ájìl} & `\textit{monter} [go up]' \\
		DOWN & \textit{sə́} & `\textit{descendre} [descend]' \\
		\lspbottomrule
	\end{tabular}
\end{table}

In the absence of a lexicon, the verbs in Table \ref{tab:somrV} are aggregated from \citet{jungraithmayrPresentationDunConte1978}, \citet{barreteauVerbeSibine1982} and \citet{millerInformationStructureSoumraye2020}. Note that the verbs \textit{há} `go' and \textit{háde} `come' have similar forms but the syllable \textit{de} is not analyzed as a suffix. This may be evidence of an erstwhile ventive extension.

\subparagraph{Directional extensions}

No motion-related verbal morphology is described for Somrai. 

\subparagraph{Caused accompanied motion (CAM)}

No information available. 

\subparagraph{Buy-sell verbs}

\citet[61]{millerInformationStructureSoumraye2020} mentions the verb \textit{kələ̂} `buy' but it is not clear how the notion `sell' is expressed. 

\subsection{East Chadic A2 (1/6)}

Only Lele is included from the six East Chadic A2 languages. No grammatical description was found for Nancere (nanc1253), Gabri (gabr1253), Kimre (kimr1241) or Kabalai (kaba1292). In a short description of Tobanga (toba1280) grammar, \citet[141]{caprileNotesLinguistiquesTobanga1978} includes a suffix \textit{-ga} called ``\textit{accompli directionnel} [directional completive/perfective]'' in an inventory of Tobanga verbal morphology. However, the meaning of the suffix is not discussed and the examples of verbs in this form do not include translations. The suffix appears several times in a short text, but not with any clear motion-related function. With such limited information, Tobanga was not included in this survey. 

\paragraph*{Lele (lele1276)} \label{lele}

A rough, but extensive sketch of Lele grammar is found in \citet{frajzyngierGrammarLele2001}. Previous descriptions of Lele grammar include short articles \citep{simonscopeNeBeMarking1982,simonscopePluralLele1993,simonscopeCommentsNominalizationLele1986} and a lexicon \citep{146082}. 

%Most basic directional meanings are lexicalized in Lele, but there is also a ventive marker and an itive adverbial. Directional meaning with manner of motion verbs is expressed in a multiverb construction. Multiverb constructions can also express (ventive) prior associated motion. Caused accompanied motion is primarily expressed using the associative preposition, though not necessarily with an over complement. There is one buy-sell verb which can be specified to mean `sell' when combined with an itive adverbial. 

%check page numbers from PDF against printed copy of Fraj

\subparagraph{Directional verbs}~

\begin{table}[H]
	\caption{Lele directional verbs}
	\label{tab:leleV}
	\begin{tabular}{lll} % add l for every additional column or remove as necessary
		\lsptoprule
		DIR & verb & gloss   \\ %table header
		\midrule
		ITIVE & \textit{è} & `go' \\
		VENT & \textit{èjè} & `come' \\
		DOWN & \textit{jìrmí} & `descend' \\
		INTO& \textit{do'b} & `enter' \\
		OUT & \textit{an} & `go/come out' \\
		\lspbottomrule
	\end{tabular}
\end{table}

Lele has intrinsic ventive and itive motion verbs, \textit{è} `go' and \textit{èjè} `come'. Most forms of the ventive form are transparently derived from the itive verb and the ventive directional marker. However, in the imperative forms, the two verbs differ: \textit{ìrà} `go.\textsc{imp}', \textit{àjè} `come.\textsc{imp}' \citep[50]{frajzyngierGrammarLele2001}. Further evidence of the lexicalization of `come' is that pronominals normally occur between the verb and the ventive marker, but this does not occur with the verb \textit{èjè} `come', indicating that [jè] is not a separate morpheme in this context (\citeauthor{frajzyngierGrammarLele2001} \citeyear[202]{frajzyngierGrammarLele2001}; \citeauthor{simonscopeNeBeMarking1982} \citeyear[228]{simonscopeNeBeMarking1982}; \citeauthor{146082} \citeyear[21]{146082}).\footnote{Despite this, Frajzyngier consistently glosses the verb \textit{èjè} `come' as two separate morphemes.}

There are verbs for inward motion, \textit{do'b} `enter', and outward motion, \textit{an} `go{\slash}come out'. There is a verb for downward motion \textit{jìrmí} `descend', but an adverbial is used to express upward motion \citep{146082,frajzyngierGrammarLele2001}.

\subparagraph{Directional extensions}~

\begin{table}[H]
	\caption{Lele directional morphology}
	\label{tab:lele}
	\begin{tabular}{llll} % add l for every additional column or remove as necessary
		\lsptoprule
		forms & source gloss & direction  & other functions  \\ %table header
		\midrule
		\textit{jè}  &  \textsc{vent}  & \textsc{vent}  &  \textsc{subs.vent} \\
		\lspbottomrule
	\end{tabular}
\end{table}

The ventive marker has a directional function when occurring with a verb of motion, as in (\ref{ex2:lele1}) and (\ref{ex2:lele2}). Note that the subject pronoun following the verb occurs between the verb and the ventive marker. \citet[196]{frajzyngierGrammarLele2001} says that the ventive marker is derived from a verb meaning `come’, but there is no evidence given for this. 

\ea\label{ex2:lele1}
\langinfo{Lele verb \textit{kìn} ‘return’ with ventive marker}{}{\cite[195]{frajzyngierGrammarLele2001}} \\
\gll kìn gé jè ná kùlbá go dà súk ni	\\
return 3\gl{pl} \gl{vent} \gl{asoc} cow \gl{dem} \gl{prep} market \gl{loc} \\
\glt `They brought a cow from the market. (The speaker is in the place to which the cow is brought.)'

\ex\label{ex2:lele2}
\langinfo{Lele verb \textit{ana} ‘exit’ with ventive marker}{}{\cite[271]{frajzyngierGrammarLele2001}} \\
\gll yàá bèì na gìdìrè na ma ná, ìrà ná kìrè kíndi kara índùwé ma gé ná ana gé jè \\
tell \gl{dat}-3\gl{m} \textsc{hypothetical} moon \textsc{hypothetical} die.\gl{imp} \gl{asoc} go.\gl{imp} \gl{asoc} road truly people human.\gl{pl} die.\gl{imp} 3\gl{pl} \gl{asoc} exit.\gl{imp} 3\gl{pl} \gl{vent}	\\
\glt `[He] told him that the moon would die and never return, and the people would die and return.'
\z

\citet[194]{frajzyngierGrammarLele2001} gives one example of a subsequent motion interpretation of a ventive marker with a non-motion main verb, shown in (\ref{ex2:lele3}). There are no other examples with an explicit subsequent motion translation, but the semantics of (\ref{ex2:lele4}) are compatible with a subsequent motion interpretation of the ventive marker. 

\ea\label{ex2:lele3}
\langinfo{Lele verb \textit{lèé} ‘eat’ with ventive marker}{}{\cite[194]{frajzyngierGrammarLele2001}} \\
\gll ŋ	lèé	jè	sìí \\
1\gl{sg}	eat	\gl{vent}	meat \\
\glt `I ate the meat [somewhere else and came here].'

\ex\label{ex2:lele4}
\langinfo{Lele verb \textit{wél} ‘sleep’ with ventive marker}{}{\cite[380]{frajzyngierGrammarLele2001}} \\
\gll tòn-iŋ	dú	na	ŋ	wél	jè	mínà	gà \\
ask-1\gl{sg}	3\gl{f}	\textsc{hypothetical}	1\gl{sg}	sleep	\gl{vent}	where	\gl{q} \\
\glt `She asked where I slept.'
\z 

In other cases, as in (\ref{ex2:lele5}) and (\ref{ex2:lele6}), the translation and context leave unclear what the function of the ventive marker is. 

\ea\label{ex2:lele5}
\langinfo{Lele verb \textit{ne} ‘make’ with ventive marker}{}{\cite[196]{frajzyngierGrammarLele2001}} \\
\gll gi	ne	jè	gùyé	bé	wéy	gà \\
2\gl{m}	make	\gl{vent}	work	\gl{dat}	who	\gl{q} \\
\glt `For whom did you work?'

\ex\label{ex2:lele6}
\langinfo{Lele verb \textit{gírbí} ‘forget’ with ventive marker}{}{\cite[384]{frajzyngierGrammarLele2001}} \\
\gll ŋ	sèn	go	gi	gírbí	jè	kòjò	kò-m	bé	kolo	ŋ	gol \\
1\gl{sg}	know	??	2\gl{m}	forget	\gl{vent}	hoe	2\gl{poss}-2\gl{m}	\gl{dat}	reason	1\gl{sg}	see \\
\glt `I know that you forgot your hoe, because I saw [that you forgot your hoe].'
\z 

The locative adverbial \textit{cáaní} ‘in/at the bush’, as in (\ref{ex2:lele7}), also has an itive function, as in (\ref{ex2:lele8}) and (\ref{ex2:lele9}). The adverbial \textit{cáaní} can occur following a direct object, unlike the ventive marker which occurs before a direct object, as in (\ref{ex2:lele5}) and (\ref{ex2:lele6}). 

\ea\label{ex2:lele7}
\langinfo{Lele (locative) adverbial \textit{cáaní} ‘in the bush’ }{}{\cite[211]{frajzyngierGrammarLele2001}} \\
\gll màdú	nè	ná	cáaní	lay \\
vitex	\textsc{copula}	\gl{asoc}	in.the.bush	also \\
\glt `The fruits of vitex are in the bush.'

\ex\label{ex2:lele8}
\langinfo{Lele verb \textit{húm} ‘gather’ with (itive) adverbial \textit{cáaní}}{}{\cite[197]{frajzyngierGrammarLele2001}} \\
\gll húm	ká∫yà	sóre	kò-m	cáaní	ɗá	bé-ŋ	kírè \\
gather.\gl{imp}	things	play	2\gl{poss}-2\gl{m}	in.the.bush	let	\gl{dat}-1\gl{sg}	road \\
\glt `Gather your toys away from here and let me go.'

\ex\label{ex2:lele9}
\langinfo{Lele verb \textit{sáde} ‘clean’ with (itive) adverbial \textit{cáaní}}{}{\cite[198]{frajzyngierGrammarLele2001}} \\
\gll ŋ	sáde	kùbàrò	cáaní \\
1\gl{sg}	clean.\gl{fut}	blood	in.the.bush \\
\glt `I will clean the blood away.'
\z 

Although \textit{cáaní} is not part of the verbal complex, it plays a role in describing buying-selling events (see below). 

\subparagraph{Caused accompanied motion (CAM)}

The primary strategy for expressing caused accompanied motion is the use of an associative preposition with a verb of motion, as in (\ref{ex2:lele16}). Note that the associative preposition can be used without an overt object, as in (\ref{ex2:lele17}). 

\ea\label{ex2:lele16}
\langinfo{Lele caused accompanied motion}{}{\cite[203]{frajzyngierGrammarLele2001}} \\
\gll èje dí ná kama \\
come 3\gl{sg} \gl{asoc} water \\
\glt `He brought water.'

\ex\label{ex2:lele17}
\langinfo{Lele caused accompanied motion}{}{\cite[154]{frajzyngierGrammarLele2001}} \\
\gll ng èjé ná túgú \\
1\gl{sg} come \gl{asoc} village \\
\glt `I bring them home.'
\z 

\subparagraph{Buy-sell verbs}

The verb \textit{kil} is a general verb of trade or exchanging goods for money. According to the context, it can be translated as either `buy', as in (\ref{ex2:lele10}) or `sell', as in (\ref{ex2:lele11}).

\ea\label{ex2:lele10}
\langinfo{Lele verb \textit{kil} translated ‘buy'}{}{\cite[163]{frajzyngierGrammarLele2001}} \\
\gll gìlkínín dàgè ná kiya kil gé kùlbá \\
Gilkinin 3\gl{pl} \gl{asoc} Kiya buy/sell 3\gl{pl} cow \\
\glt `Gilkinin and Kiya bought a cow.'

\ex\label{ex2:lele11}
\langinfo{Lele verb \textit{kil} translated ‘sell'}{}{\cite[120]{frajzyngierGrammarLele2001}} \\
\gll kire kil gùjù ná tòngú kè-y lìŋdà \\
Kire buy/sell millet \gl{asoc} beans 2\gl{poss}-3\gl{m} yesterday \\
\glt `Kire sold his millet and beans yesterday.'

\z 

The `sell' interpretation of \textit{kil} can be reinforced by the adverbial \textit{cáaní} `in the bush' in its itive function, as in (\ref{ex2:lele12}).

\ea\label{ex2:lele12}
\langinfo{Lele verb \textit{kil} meaning ‘sell' with itive adverbial \textit{cáaní}}{}{\cite[161]{frajzyngierGrammarLele2001}} \\
\gll kil	gé	kàmyá	tùmò	cáaní \\
buy/sell	3\gl{pl}	milk	before	in.the.bush \\
\glt `They used to sell milk.'
\z

\subsection{East Chadic A3 (2/2)}

\paragraph*{Kera (kera1255)} \label{kera}

Information on Kera grammar is from a reference grammar \citep{ebertSpracheUndTradition1976}. 

\subparagraph{Directional verbs}~

\begin{table}[H]
	\caption{Kera directional verbs} 
	\label{tab:keraV}
	\begin{tabular}{lll} % add l for every additional column or remove as necessary
		\lsptoprule
		DIR & verb & gloss   \\ %table header
		\midrule
		ITIVE & \textit{ɗé} & `\textit{gehen}, \textit{aller} [go]' \\
		VENT & \textit{bi} & `\textit{kommen}, \textit{venir} [come]' \\
		INTO & \textit{kélé} & `\textit{eintreten}, \textit{entrer} [enter]'\\
%		INTO/OUT & \textit{kél wə́ra} & `\textit{hinausgehen}, \textit{sortir} [exit]'\\
		UP/DOWN & \textit{lí} & `(\textit{be})\textit{steigen}, \textit{monter} [go up]' \\
%		UP/DOWN & \textit{lí wə́ra} & `(\textit{hinaub})\textit{steigen}, \textit{descendre} [go down]' \\
		\lspbottomrule
	\end{tabular}
\end{table}

The ``totality'' extension \textit{wə́ra} is said to express the completion of an action, especially when the direct object no longer exists or is no longer present as a result of the action \citep[111]{ebertSpracheUndTradition1979}. It is used in the lexicon to show the outward meaning of \textit{kélé}, as in \textit{kél wə́ra} `exit', and also to show the downward meaning of \textit{lí}, as in \textit{lí wə́ra} `go down' \citep{ebertSpracheUndTradition1976}. The verb \textit{lí} is glossed as upward motion, but see example (\ref{ex2:kera2}) where it is interpreted as downward motion without the totality extension.

\subparagraph{Directional extensions}~

\begin{table}[H]
	\caption{Kera directional morphology}
	\label{tab:kera}
	\begin{tabular}{llll} % add l for every additional column or remove as necessary
		\lsptoprule
		forms & source gloss & direction  & other functions  \\ %table header
		\midrule
		\textit{-dà}  &  \textsc{da}  & \textsc{vent}   &  \textsc{subs.vent} \\
		\textit{-na}  &  \textsc{dist}  & \textsc{itive}   &  \\
		\lspbottomrule
	\end{tabular}
\end{table}

\citet[113--116]{ebertSpracheUndTradition1979} describes several functions of the marker \textit{-dà}. When \textit{-dà} follows a motion verb, it has a ventive directional function: ``The presence and absence of -dà corresponds to the difference between coming and going in verbs of movement'' \citep[114]{ebertSpracheUndTradition1979}. With intransitive motion verbs, as in (\ref{ex2:kera1}) and (\ref{ex2:kera2}), \textit{-dà} has a ventive directional meaning. 

\ea\label{ex2:kera1}
\langinfo{Kera verb \textit{kə́láŋ} ‘enter' with ventive marker}{}{\cite[113]{ebertSpracheUndTradition1979}} \\
\gll kə́láŋ-dà	\\
enter-\gl{vent}	\\
\glt `came in' [German: `\textit{kam herein}']

\ex\label{ex2:kera2}
\langinfo{Kera verb \textit{lúŋ} ‘ascend/descend' with ventive marker}{}{\cite[114]{ebertSpracheUndTradition1979}} \\
\gll Hùlùm bə̀ Péve, wə lúŋ-dà, wə lúŋ Áw-dà, Dòrè-dà.	\\
person ?? Peve ?? ascend/descend-\gl{vent} ?? ascend/descend Aw-\gl{vent} Dore-\gl{vent}	\\
\glt `The Peve, they came down, they came down from Aw, from Dore.' \newline [German: `\textit{Die Peve, sie kamen herab, sie kamen von Aw herunter, von Dore.}]	'
\z 

The ventive marker can also occur (redundantly) with the ventive motion verb \textit{bi} `come', as in (\ref{ex2:kera3}), but not with the itive motion verb \textit{ɗé} `go'.

\ea\label{ex2:kera3}
\langinfo{Kera verb \textit{bi} ‘come' with ventive marker}{}{\cite[113]{ebertSpracheUndTradition1979}} \\
\gll bàŋ-dà	\\
come-\gl{vent}	\\
\glt `came here' [German: `\textit{kam her}']
\z 	

Examples of \textit{-dà} with non-motion verbs have a subsequent caused accompanied motion interpretation, as in (\ref{ex2:kera4}). 

\ea\label{ex2:kera4}
\langinfo{Kera verb \textit{tə́káŋ} ‘reap' with ventive marker}{}{\cite[114]{ebertSpracheUndTradition1979}} \\
\gll tə́káŋ-dà	\\
reap-\gl{vent}	\\
\glt `reaped (and brought) here' [German: `\textit{erntete} (\textit{und brachte}) \textit{her}']
\z 

The verbal morpheme \textit{wə́ra}, described as an aspectual marker, perhaps similar to what are called totality markers in other Chadic languages, can appear between the verb root and the ventive marker. The totality marker has some idiosyncratic interpretations, such as in (\ref{ex2:kera8}) where, in combination with the verb \textit{kə́láŋ}, elsewhere translated `enter', it rather means `exit'. In combination with the ventive marker \textit{-dà}, as in (\ref{ex2:kera9}), the exiting motion is described as occurring in a direction toward the deictic center.\footnote{It is not immediately clear why there is consistently a space in the transcription between the verb root and totality marker, but a space is found in the same place in the Kera orthography (Jackie Hainaut, personal communication).}

\ea\label{ex2:kera8}
\langinfo{Kera verb \textit{kə́láŋ} ‘enter' with totality marker}{}{\cite[114]{ebertSpracheUndTradition1979}} \\
\gll kə́láŋ wə́ra	\\
enter \gl{tot}	\\
\glt `went out' [German: `\textit{ging hinaus}']

\ex\label{ex2:kera9}
\langinfo{Kera verb \textit{kə́láŋ} ‘enter' with totality and ventive markers}{}{\cite[114]{ebertSpracheUndTradition1979}} \\
\gll kə́láŋ wár-dà	\\
enter \gl{tot}-\gl{vent}	\\
\glt `came out' [German: `\textit{kam heraus}']
\z 

The marker \textit{-dà} can also follow locative nouns, as seen in (\ref{ex2:kera2}). In this context, \textit{-dà} appears to mark source arguments. However, \citet[114]{ebertSpracheUndTradition1979} states that this is not always the case, and further argues that the function of \textit{-dà} following locatives need not be distinguished from its ventive directional function following a verb of motion.  

The morpheme \textit{-ná} is described as a counterpart of \textit{-dà} but with distal meaning \citep[115]{ebertSpracheUndTradition1979}. In the few examples Ebert gives, the distal marker is either co-occurring with the totality marker (see (\ref{ex2:kera10}) below) or appears to have a postpositional function. However, through personal communication with Jackie Hainaut, it was confirmed that in at least some contexts it is possible for the itive \textit{-na} marker to appear suffixed to a verb of motion with a directional function, as in (\ref{ex2:kera7}). 

\ea\label{ex2:kera7}
\langinfo{Kera verb \textit{kə́láŋ} ‘enter' with itive marker}{}{Jackie Hainaut, personal communication} \\
\gll kə́láŋ-na	\\
enter-\gl{itive}	\\
\glt `went in'
\z 

\subparagraph{Caused accompanied motion (CAM)}

Directed CAM may be expressed by the use of the ventive marker \textit{-dà} with a verb glossed `take', as in (\ref{ex2:kera5}) and (\ref{ex2:kera6}).

\ea\label{ex2:kera5}
\langinfo{Kera verb \textit{hàŋ} ‘take' with ventive marker}{}{\cite[114]{ebertSpracheUndTradition1979}} \\
\gll hàŋ-dà	\\
take-\gl{vent}	\\
\glt `he took (and brought) here' [German: `\textit{nahm} (\textit{und brachte})']

\ex\label{ex2:kera6}
\langinfo{Kera verb \textit{hàw} ‘take' with ventive marker}{}{\cite[116]{ebertSpracheUndTradition1979}} \\
\gll hàw-dà 	\\
take-\gl{vent}	\\
\glt `Bring it here!' [German: `\textit{Bring es her}!']
\z 

The itive (or distal) morpheme \textit{-na} is also found in expressions of (itive) CAM, but the few examples available always show it alongside the totality morpheme \textit{wəra}, as in (\ref{ex2:kera10}).

\ea\label{ex2:kera10}
\langinfo{Kera verb \textit{hàŋ} ‘take' with totality and itive markers}{}{\cite[116]{ebertSpracheUndTradition1979}} \\
\gll wə hùmùŋ áskáŋ war-ná ádə̀-ná	\\
3\gl{sg} take fish \gl{tot}-\gl{itive} ??-\gl{itive}	\\
\glt `He took the fish away (and brought them) there.' \newline [German: `\textit{Er nahm die Fische weg} (\textit{und brachte sie}) \textit{dorthin}.']
\z 

\subparagraph{Buy-sell verbs}

The verb \textit{gùsí} `buy' expresses the notion `sell' when followed by the ``totality'' marker \textit{wə́ra}, as in \textit{gùsí wə̀ra} `sell' \citep[53]{ebertSpracheUndTradition1976}. 

\paragraph*{Kwang (kwan1285)} \label{kwan}

Information on the grammar of Kwang is from a grammar sketch published as a special issue of a journal \citep{lenssenStudienVerbIm1984}. 

\largerpage
\subparagraph{Directional verbs}~

\begin{table}[H]
	\caption{Kwang directional verbs} 
	\label{tab:kwanV}
	\begin{tabular}{llll} % add l for every additional column or remove as necessary
		\lsptoprule
		DIR & verb & gloss &  \\ %table header
		\midrule
		ITIVE & \textit{ɗe} & `\textit{aller} [go]' & \\
		VENT & \textit{ɗéːna} & `\textit{venir} [come]' & \\
		INTO& \textit{gálkájaŋ} & `\textit{entrer} [enter]' & Mobu lect \\
		INTO& \textit{báyé} & `\textit{entrer} [enter]' & Ngam lect \\
		OUT & \textit{gale} & `\textit{sortir} [exit]' & Mobu lect \\
		OUT & \textit{gáljóŋo} & `\textit{sortir} [exit]' & Ngam lect \\
		UP & \textit{ale} & `\textit{monter} [go up]' &  \\
		DOWN & \textit{jírbé} & `\textit{descendre} [descend]' & Mobu lect \\
		DOWN & \textit{jìrpe, jàːpe} & `\textit{descendre} [descend]' & Ngam lect \\
		\lspbottomrule
	\end{tabular}
\end{table}

\subparagraph{Directional extensions}

The description of Kwang by \citet{lenssenStudienVerbIm1984} does not include any (productive) motion-related verbal morphology. It is noted that the verb \textit{ɗéːna} `come' is transparently related to \textit{ɗe} `go'. However, this is described as a lexicalized morpheme and not a regular ventive suffix in the verbal morphology \citep[7, 55]{lenssenStudienVerbIm1984}. 

Another sign of possible erstwhile directional morphology is found on the ends of the verbs \textit{gálkájaŋ} `\textit{entrer} [enter]' and  \textit{gáljóŋo} `\textit{sortir} [exit]' in the Ngam dialect when compared with the shorter verb root form in the Mobu dialect verb: \textit{gale} `\textit{sortir} [exit]' \citep[55]{lenssenStudienVerbIm1984}. 

\subparagraph{Caused accompanied motion (CAM)}

No information available. 

\subparagraph{Buy-sell verbs}

The lexical entry for the verb \textit{kèle} `buy' (Ngam) contains the sub-entry \textit{kèlè sə́ŋ} `sell' \citep[67--68]{lenssenStudienVerbIm1984}. It is not clear whether \textit{sə́ŋ} has any other function in the language.

\section{East Chadic B (Guera) (11/22)}

Grammatical descriptions are available for half of the East Chadic B languages. There is evidence of directional extensions in only one of these 11 East Chadic B languages. 

%Correspondingly, each language has a set of lexical motion verbs for specific types of directional motion, as seen in Table \ref{tab:EastBV}. Limited information is available about the expression of directed CAM. Four languages allow a valency-increasing alternation with a directional motion verb to express directed CAM. In the case of Barayin, the ventive and itive CAM verbs are now lexicalized, but they are transparently (though not predictably) derived historically from directional motion verbs that were modified by a valency-increasing suffix. Where specific forms have developed to distinguish the meaning SELL from BUY (or TRADE), this is normally done with a separate verb root. The sources of these roots are not always clear, and their forms are not transparently cognate. In one case, Migaama, the SELL meaning can be specified by the use of an oblique suffix on the verb stem meaning TRADE/BUY. 

%\begin{table}[H]
%	\caption{East Chadic B directional verbs}
%	\label{tab:EastBV}
%	\begin{tabular}{lllllll} % add l for every additional column or remove as necessary
%		\lsptoprule
%		language & VENT & ITIVE  & UP & DOWN & INTO & OUT  \\ %table header
%		\midrule
%		Birgit (birg1239) & \textit{'àsí} & \textit{kàati} & \textit{tàadí} & \textit{zìigí} & \textit{ònyí} & \textit{kàati} \\	
%		Bidiya (bidi1241) & \textit{ás}  & \textit{ɗàwàn} & \textit{àcal} & \textit{beeg, zìil} & \textit{ompu} &  \textit{amàl}  \\	
%		Dangla (dang1276) & \textit{ase} & \textit{ɓaawe} & \textit{coone} & \textit{paay} &  \textit{un̰je} & \textit{amile}  \\	
%		Mogum (mogu1251) & \textit{asa} & \textit{daŋi}  & \textit{yapi} & \textit{beri} & \textit{uɲi, oɲa} & \textit{kati} \\	
%		Migaama (miga1249) & \textit{ásáw}  & \textit{áɗ-} & \textit{kálpó} & \textit{bèerò} & \textit{ùnnyò} & \textit{òomò} \\	
%		Mubi (mubi1246) & \textit{ságé} & \textit{ɓów} & \textit{nàagé} & \textit{gìidí} & \textit{rúuɗí} & \textit{kàadé} \\	
%		Mukulu (muku1242) & \textit{éttó} & \textit{èngílè} & \textit{kìlí} & \textit{ɗéègé, zábbè} & \textit{díinè} & \textit{áɗɗè} \\	
%		Mawa (mawa1270) & \textit{ooboŋ}  &\textit{ɓiaaŋ} & \textit{teedeŋ}  & \textit{jaaye}  & \textit{koyoŋ} &  \textit{giñiŋ} \\	
%		Sokoro (soko1263) &  \textit{obiŋ} & \textit{ɓeŋ} & \textit{təɗi} & \textit{wajir} & \textit{koypa} & \textit{gin̰}  \\	
%		Barayin (bare1279) & \textit{ājó, síi} & \textit{kóló} &  \textit{tado}  & \textit{jàŋgó} & \textit{tōpó} &  \textit{gùsò} \\	
%		\lspbottomrule
%	\end{tabular}
%\end{table}

%\begin{table}[H]
%	\caption{East Chadic B BUY-SELL verbs}
%	\label{tab:EastBBUY}
%	\begin{tabular}{lll} % add l for every additional column or remove as necessary
%		\lsptoprule
%		language & TRADE/BUY & SELL    \\ %table header
%		\midrule
%		Birgit (birg1239) & \textit{ḍòoyí}  &  \textit{sùumí}  \\	
%		Bidiya (bidi1241) &  \textit{gìdày} &  \textit{gìdày (ùnda)} \\
%		Dangla (dang1276) & \textit{gidiye}, \textit{gùsí} & \\	
%		Mogum (mogu1251) & \textit{ɗawi} & \textit{daga}  \\	
%		Migaama (miga1249) & \textit{gíɗáw}  & \textit{geɗek-ki-ko} \\	
%		Mubi (mubi1246) & \textit{ôjé} & \\	
%		Mukulu (muku1242) & \textit{ôwilâ} & \textit{ɗáàbé}  \\	
%		Mawa (mawa1270) & \textit{ebereŋ} &  \\	
%		Sokoro (soko1263) & \textit{epir} & \textit{diging}   \\	
%		Barayin (bare1279) & \textit{gōrō} &  \textit{bìtó} \\
%		\lspbottomrule
%	\end{tabular}
%\end{table}

\subsection{East Chadic B1 (6/12)}

Half of the 12 East Chadic B1 languages are included here. Kajakse (kaja1254), Masmaje (masm1239) and Toram (tora1267) are described by \citet{alioPreliminariesEtudeLangue2004} but not in enough depth to be included. The description of Zerenkel (zire1244) is a work in progress, but the relevant data is not available at the moment. Jonkor Bourmataguil (jonk1238) and Mabire (mabi1242) have essentially no documentation beyond short word lists, and the latter is now considered dormant. 

\paragraph*{Birgit (birg1239)}

Information on the grammar of Birgit is from a grammar sketch published as a journal article \citep{jungraithmayrBirgitOsttschadischeSprache2004}. 

\newpage
\subparagraph{Directional verbs}~

\begin{table}[H]
	\caption{Birgit directional verbs} 
	\label{tab:birgV}
	\begin{tabular}{lll} % add l for every additional column or remove as necessary
		\lsptoprule
		DIR & verb & gloss   \\ %table header
		\midrule
		ITIVE/OUT & \textit{kàati} & `\textit{aller}, \textit{sortir} [go, go out]' \\
		VENT & \textit{'àsí} & `\textit{venir} [come]' \\
		INTO& \textit{ònyí} & `\textit{entrer} [enter]' \\
		UP & \textit{tàadí} & `\textit{monter} [ascend]' \\
		DOWN & \textit{zìigí} & `\textit{descendre} [descend]' \\
		\lspbottomrule
	\end{tabular}
\end{table}

While the verb \textit{kàati} is glossed as `\textit{aller}, \textit{sortir} [go, go out]' in the lexicon, the limited amount of Birgit data do not include any examples that show this verb being used as an outward motion verb. It is possible that the `\textit{sortir}' gloss rather relates to the sense of leaving often associated with itive motion verbs. In the absence of further evidence, the gloss given will be taken at face value as an outward directional meaning. 

\subparagraph{Directional extensions}~

The short grammar sketch does not include any motion-related verbal morphology. 

\subparagraph{Caused accompanied motion (CAM)}

It appears that ventive CAM may be expressed by a form derived from the motion verb \textit{àsí} `come' followed by a causative or transitivizing marker: \textit{àsí dò} `\textit{ammener} [bring]' \citep[351]{jungraithmayrBirgitOsttschadischeSprache2004}. The marker \textit{dò} is found in a few other lexical entries for transitive verbs and in some cases an intransitive/non-causative version of the verb is also given, e.g., \textit{àllàmi} `\textit{apprendre} [learn]' and \textit{àllàmi dò} `\textit{enseigner} [teach]' \citep[350]{jungraithmayrBirgitOsttschadischeSprache2004}.

\subparagraph{Buy-sell verbs}

There are two separate verbs in the Birgit lexicon: \textit{ḍòoyí} `buy' and \textit{sùumí} `sell', the latter of which is said to be from Arabic.

\newpage
\paragraph*{Bidiya (bidi1241)} \label{bidi}

\subparagraph{Directional verbs}~

\begin{table}[H]
	\caption{Bidiya directional verbs} 
	\label{tab:bidiV}
	\begin{tabular}{lll} % add l for every additional column or remove as necessary
		\lsptoprule
		DIR & verb & gloss   \\ %table header
		\midrule
		ITIVE & \textit{ɗàwàn} & `\textit{aller} [go]' \\
		VENT & \textit{ás} & `\textit{venir} [come]' \\
		INTO& \textit{ompu} & `\textit{entrer} [enter]' \\
		OUT & \textit{amàl} & `\textit{sortir} [exit]' \\
		UP & \textit{àcal} & `\textit{monter} [go up]' \\
		DOWN & \textit{beeg, zìil} & `\textit{descendre} [go down]' \\
		\lspbottomrule
	\end{tabular}
\end{table}

Verbs in Table \ref{tab:bidiV} are from \citet{alioLexiqueBidiya1989}.

\subparagraph{Directional extensions}
\citet{alioEssaiDescriptionLangue1986} does not describe any motion-related verbal morphology in Bidiya. 

\subparagraph{Caused accompanied motion (CAM)}

The verb \textit{èe} `\textit{apporter} [carry]' can express caused accompanied motion, but it is unclear how directed CAM is expressed in Bidiya. This verb appears to be used without regard to deixis in the texts included in \citet{seibelParticipantReferenceBidiya2015}.

\subparagraph{Buy-sell verbs}

In the lexicon of \citet[78]{alioLexiqueBidiya1989} the notion `sell' is shown to be derived from the verb \textit{gìdày} `buy' when modified by an adverb \textit{ùnda} `\textit{devant} [in front of]' as in \textit{gìdày ùnda} `sell'. However, \citet[29, 44, 77]{seibelParticipantReferenceBidiya2015} includes examples where \textit{gìdày} is translated as `sell' without the adverbial modifier.

\paragraph*{Dangla (dang1276)} \label{dang}

There are three identified subgroups of Dangla: West, Central and East. \citet{shayGrammarEastDangla1999} describes East Dangla, and \citet{peustSupplementsWestDangla2016} describes West Dangla based on materials collected by Jacques Fédry in 1966 and 1967 \citep{fedryDictionnaireDangaleatTchad1971}.  There is also an online dictionary of Dangla \citep{webonaryDictionnaireDangaleat2016}.

\newpage
\subparagraph{Directional verbs}~

\begin{table}[H]
	\caption{Dangla directional verbs \citep{webonaryDictionnaireDangaleat2016}} 
	\label{tab:danglaV}
	\begin{tabular}{lll} % add l for every additional column or remove as necessary
		\lsptoprule
		DIR & verb & gloss   \\ %table header
		\midrule
		ITIVE & \textit{ɓaawe} & `go' \\
		VENT & \textit{ase} & `come' \\
		UP & \textit{coone} & `ascend' \\
		DOWN & \textit{paay} & `descend' \\
		INTO& \textit{un̰je} & `enter' \\
		OUT & \textit{amile} & `exit' \\

		\lspbottomrule
	\end{tabular}
\end{table}

\subparagraph{Directional extensions}~

Dangla does not have directional verbal morphology. 

\subparagraph{Caused accompanied motion (CAM)}

Caused accompanied motion is primarily expressed by the verb \textit{iye} `carry' (French: `\textit{apporter}') \citep{webonaryDictionnaireDangaleat2016}. This verb does not have any inherent directional meaning, and none of the available sources indicate any grammaticalized means of specifying the direction of the verb. 

\subparagraph{Buy-sell verbs}

Dangla has a verb that can be translated as either `buy' or `sell': \textit{gidiye} in East Dangla  \citep{shayGrammarEastDangla1999} and West Dangla \citep{peustSupplementsWestDangla2016} and \textit{gùsí} in Central Dangla \citep[32]{burkeIntroductionVerbalSystem1995}. The online Dangla dictionary also reports the form \textit{suugiye} \citep{webonaryDictionnaireDangaleat2016}. The only source that gives multiple examples of this verb in context is \citet{shayGrammarEastDangla1999}. There are no clear patterns of any grammatical means of specifying a `buy' or `sell' interpretation of the verb, other than the general context.  

\paragraph*{Mogum (mogu1251)}

Information on Mogum is primarily from unpublished notes created during workshops run by SIL for the purpose of helping Mogum speakers develop a writing system for their language. 

\newpage
\subparagraph{Directional verbs}~

\begin{table}[H]
	\caption{Mogum directional verbs} 
	\label{tab:moguV}
	\begin{tabular}{lll} % add l for every additional column or remove as necessary
		\lsptoprule
		DIR & verb & gloss   \\ %table header
		\midrule
		ITIVE & \textit{daŋi} & `\textit{aller} [go]' \\
		VENT & \textit{asa} & `\textit{venir} [come]' \\
		UP & \textit{yapi} & `\textit{monter} [ascend]' \\
		DOWN & \textit{beri} & `\textit{descendre} [descend]' \\
		INTO& \textit{uɲi, oɲa} & `\textit{entrer} [enter]' \\
		OUT & \textit{kati} & `\textit{sortir} [exit]' \\
		\lspbottomrule
	\end{tabular}
\end{table}

\subparagraph{Directional extensions}

There is no information indicating that Mogum has any motion-related verbal morphology.

\subparagraph{Caused accompanied motion (CAM)}

From a handful of examples, it appears that directed CAM may be expressed by a transitive form of the ventive motion verb. Compare \textit{No asa} `\textit{Je viens} [I come]' with \textit{No asaa} `\textit{J'amène} [I bring]'. The status of the second \textit{a} at the end of the verb is unclear, as this is not a pattern with causatives. In another example, a direct object suffix is included: \textit{No asaak} `\textit{Je les amène} [I bring them]'. 

\subparagraph{Buy-sell verbs}

There are two verbs in Mogum: \textit{ɗawi} `buy' and \textit{daga} `sell'. 

\paragraph*{Migaama (miga1249)} \label{miga}

A short sketch of Migaama grammar is included in a lexicon of the language \citep{jungraithmayrLexiqueMigamaMigamafrancais1992} and there is also an unpublished grammar sketch done by Migaama speakers in an SIL workshop \citep{moussaDTLNovembre2010n.d.}.

\newpage
\subparagraph{Directional verbs}~

\begin{table}[H]
	\caption{Migaama directional verbs} 
	\label{tab:migaV}
	\begin{tabular}{lll} % add l for every additional column or remove as necessary
		\lsptoprule
		DIR & verb & gloss   \\ %table header
		\midrule
		ITIVE & \textit{áɗ-} & `\textit{aller} [go]' \\
		VENT & \textit{ásáw} & `\textit{venir}, \textit{arriver} [come, arrive]' \\
		UP & \textit{kálpó} & `\textit{monter} [go up]' \\
		DOWN & \textit{bèerò} & `\textit{descendre} [go down]' \\
		INTO& \textit{ùnnyò} & `\textit{entrer} [enter]' \\
		OUT & \textit{òomò} & `\textit{sortir} [exit]' \\
		\lspbottomrule
	\end{tabular}
\end{table}

\subparagraph{Directional extensions}~

Migaama does not have directional verbal morphology. 

\subparagraph{Caused accompanied motion (CAM)}

There is a general verb of caused accompanied motion (CAM) \textit{ítìrò} `\textit{rapporter}, \textit{amener} [bring]' but directed CAM in Migaama can be expressed using a suffix \textit{-ki(ki)} labeled `causative' by \citet{ramatMorphologieVerbeMigaama2003} and `oblique' by \citet{robertsObliquesMawa2016}. When combined with a directional motion verb, the interpretation is directed CAM, as in (\ref{ex2:miga3}) and (\ref{ex2:miga4}).

\ea\label{ex2:miga3}
\langinfo{Migaama verb \textit{ásáw} `come' with oblique suffix}{}{\cite[70]{ramatMorphologieVerbeMigaama2003}} \\
\gll  Ki askiti aɗe dombo? \\
2\gl{sg} come-\gl{obl} \gl{q} hoe \\
\glt `Did you bring the hoe?' [French: `\textit{As-tu apporté la houe?}']

\ex\label{ex2:miga4}
\langinfo{Migaama verb \textit{kaato} `go' with oblique suffix}{}{\cite[71]{ramatMorphologieVerbeMigaama2003}} \\
\gll káàtì-(kì)kì-kô \\
go-\gl{obl}-\gl{obj}.3\gl{pl} \\
\glt `brought them' [French: `\textit{les a emportés}']
\z

With other verbs, this suffix has interpretations relating to instrument or accompaniment, as in (\ref{ex2:miga5}) and (\ref{ex2:miga6}).

\ea\label{ex2:miga5}
\langinfo{Migaama verb \textit{dew} `kill' with oblique suffix}{}{\cite[71]{ramatMorphologieVerbeMigaama2003}} \\
\gll Giimi dew-(ki)ki-ti bunduk gaaruwi. \\
people tuer-\gl{obl}-\gl{obj}.3\gl{sg}.\gl{f} rifle animal \\
\glt `People kill animals with the help of a rifle.' \newline [French: `\textit{Les gens tuent les animaux á l'aide du fusil.}']

\ex\label{ex2:miga6}
\langinfo{Migaama verb \textit{tíyáw} `eat' with oblique suffix}{}{\cite[71]{ramatMorphologieVerbeMigaama2003}} \\
\gll téè-(kì)kî \\
eat-\gl{obl}\\
\glt `ate with him' [French: `\textit{a mangé avec lui}']
\z 

\subparagraph{Buy-sell verbs}

The verb \textit{gíɗáw} can be translated `buy' or `sell' according to the context. The `sell' interpretation can be enforced by using the oblique suffix, as in (\ref{ex2:miga7}).

\ea\label{ex2:miga7}
\langinfo{Migaama verb \textit{gíɗáw} `buy/sell' with oblique suffix}{}{Sakine Ramat, personal communication} \\
\gll Na  geɗek-ki-ko iini \\
1\gl{sg}  sell-\gl{obl}-\gl{obj}.3\gl{pl} goat \\
\glt `I sell goats.'
\z 

\paragraph*{Mubi (mubi1246)}

There have been several descriptions of Mubi including an MA thesis on the verbal morphology  \citep{prickettPhonologyMorphologyVerb2012}, a grammar sketch and lexicon \citep{jungraithmayrLangueMubiKaan2013} and a journal article \citep{peustProgressMubiStudies2014}. 

\subparagraph{Directional verbs}~

\begin{table}[H]
	\caption{Mubi directional verbs} 
	\label{tab:mubiV}
	\begin{tabular}{lll} % add l for every additional column or remove as necessary
		\lsptoprule
		DIR & verb & gloss   \\ %table header
		\midrule
		ITIVE & \textit{ɓów} & `\textit{aller} [go]' \\
		VENT & \textit{ságé} & `\textit{venir} [come]' \\
		UP  & \textit{nàagé} & `\textit{monter} [go up]' \\
		DOWN & \textit{gìidí} & `\textit{descendre} [go down]' \\ 
		INTO& \textit{rúuɗí} & `\textit{entrer} [enter]' \\
		OUT & \textit{kàadé} & `\textit{sortir} [exit]' \\
		\lspbottomrule
	\end{tabular}
\end{table}

\subparagraph{Directional extensions}

There is no evidence of directional verbal extensions in Mubi. 

\subparagraph{Caused accompanied motion (CAM)}

Ventive caused accompanied motion can be expressed by the causative form of the ventive directional verb (sometimes differing only in tone from the non-causative form): \textit{sàgé} \citep[85, 95]{jungraithmayrLangueMubiKaan2013}. \citet[118]{prickettPhonologyMorphologyVerb2012} also lists the verb \textit{ɟajaw} `bring'. 

\subparagraph{Buy-sell verbs}

Buying and selling events are expressed by the same verb:  \textit{ôjé} `\textit{acheter}, \textit{vendre} ; \textit{faire du commerce} [sell, buy, trade]' \citep[192]{jungraithmayrLangueMubiKaan2013}. Any grammatical means of distinguishing the two meanings is not discussed. 

\subsection{East Chadic B2 (1/1)}

\paragraph*{Mukulu (muku1242)}

\subparagraph{Directional verbs}~

\begin{table}[H]
	\caption{Mukulu directional verbs} 
	\label{tab:mukuV}
	\begin{tabular}{lll} % add l for every additional column or remove as necessary
		\lsptoprule
		DIR & verb & gloss   \\ %table header
		\midrule
		ITIVE & \textit{èngílè, étté, cînycîŋè} & `\textit{aller} [go]' \\
		VENT & \textit{éttó} & `\textit{venir}, \textit{aller} [come, go]' \\
		UP & \textit{kìlí} & `\textit{monter} [ascend]' \\
		DOWN & \textit{ɗéègé, zábbè} & `\textit{descendre} [descend]' \\
		INTO& \textit{díinè} & `\textit{entrer} [enter]' \\
		OUT & \textit{áɗɗè}  & `\textit{sortir} [exit]' \\
		\lspbottomrule
	\end{tabular}
\end{table}

The verb \textit{étté} is glossed `\textit{aller, partir; venir; amener, apporter} [go, leave, come, bring]' and the verb \textit{éttó} is glossed `\textit{venir, arriver, s’approcher, amener, apporter} [come, arrive, approach, bring]' \citep[96]{jungraithmayrLexiqueMokilko1990}. There is no indication of how these various meanings are disambiguated, and nothing in the discussion of the verbal morphology suggesting a morphological relationship between these two verbs despite their similar forms.

\subparagraph{Directional extensions}

No motion-related verbal morphology is described by \citet{jungraithmayrLexiqueMokilko1990}. 

\subparagraph{Caused accompanied motion (CAM)}

There is a general caused accompanied motion verb: \textit{ámɓè} `\textit{amener, prendre, ramasser, porter, lever, soulever} [bring, take, gather, carry, lift, lift up]' \citep[58]{jungraithmayrLexiqueMokilko1990}. Two directional motion verbs also have glosses related to CAM but without an explanation of the context in which such an interpretation is possible.  

\subparagraph{Buy-sell verbs}

Both buying and selling can be expressed by the verb \textit{ôwilâ} `buy, sell' \citep[159]{jungraithmayrLexiqueMokilko1990}. There is no explanation of the context in which each interpretation arises. Another verb also expresses selling: \textit{ɗáàbé} `sell' \citep[87]{jungraithmayrLexiqueMokilko1990}.

\subsection{East Chadic B3 (3/7)}

Three of the seven East Chadic B3 languages are included. Only minimal information exists on Tamki (tamk1242), Ubi (ubii1238), Boor (boor1242) \citep[documented for the first time by][]{robertsInitialFindingsBoor2021} and Miltu (milt1241) \citep[see the 1972 fieldnotes published as][]{jungraithmayrNotesGaliMiltu2019}.

\paragraph*{Mawa (mawa1270)} \label{mawa}

The available information on Mawa is an unfinished manu\-script from an SIL Chad workshop held in October 2012 \citep{hisseneEsquisseGrammaticaleLangue2010} and two published phonological analyses \citep{robertsPalatalizationLabializationMawa2009,robertsToneSystemMawa2012}.

\subparagraph{Directional verbs}~

\begin{table}[H]
	\caption{Mawa directional verbs} 
	\label{tab:mawaV}
	\begin{tabular}{lll} % add l for every additional column or remove as necessary
		\lsptoprule
		DIR & verb & gloss   \\ %table header
		\midrule
		ITIVE & \textit{ɓiaaŋ} & `\textit{aller} [go]' \\
		VENT & \textit{ooboŋ} & `\textit{venir} [come]' \\
		UP & \textit{teedeŋ} &  `\textit{monter} [ascend]' \\
		DOWN & \textit{jaaye} &  `\textit{descendre} [descend]' \\
		INTO& \textit{koyoŋ} &  `\textit{entrer} [enter]' \\
		OUT & \textit{giñiŋ} &  `\textit{sortir} [exit]' \\
		\lspbottomrule
	\end{tabular}
\end{table}

\subparagraph{Directional extensions}~

Mawa does not appear to have any directional verbal morphology.

\subparagraph{Caused accompanied motion (CAM)}

Caused accompanied motion is expres\-sed by adding the ``oblique'' suffix \textit{-r} to a directional verb. The verb \textit{oobiriŋ} `bring', `come with'  is derived from \textit{ooboŋ} `come' (James Roberts, personal communication). It is also possible to derive a CAM verb from the directional verb \textit{ɓiaaŋ} `go'.

\subparagraph{Buy-sell verbs}

The verb \textit{ebereŋ} only means `buy' and a different verb is used for `sell' (James Roberts, personal communication).

\paragraph*{Saba (saba1276)}

The available information on Saba is an unfinished manuscript from an SIL Chad workshop held in October 2012 \citep{abakarEsquisseGrammaticaleLanguen.d.} and the recently published fieldnotes of \citet{jungraithmayrPreliminairesEtudeSaba2020}. These sources cover verbal morphology and do not give any indication of directional verbal morphology. Information on the lexicon is very limited. 

\paragraph*{Sokoro (soko1263)} \label{soko}

Aspects of Sokoro grammar are described in an unpublished grammar sketch \citep{yarangaEsquisseGrammaticaleLangue2019}. Examples used in this section are taken from an earlier draft of this grammar sketch which mostly also occur in the final version \citep{yarangaEsquiseGrammaticaleLanguen.d.}.

\subparagraph{Directional verbs}~

\begin{table}[H]
	\caption{Sokoro directional verbs} 
	\label{tab:sokoV}
	\begin{tabular}{lll} % add l for every additional column or remove as necessary
		\lsptoprule
		DIR & verb & gloss   \\ %table header
		\midrule
		ITIVE & \textit{ɓeŋ} & `\textit{aller} [go]' \\
		VENT & \textit{ɔbiŋ} & `\textit{venir} [come]' \\
		UP & \textit{taɗiŋ} &  `\textit{monter} [ascend]' \\
		DOWN & \textit{wajir} & `\textit{descendre} [descend]' \\
		INTO& \textit{kɔypa} &  `\textit{entrer} [enter]' \\
		OUT & \textit{gin̰} & `\textit{sortir} [exit]' \\
		\lspbottomrule
	\end{tabular}
\end{table}

\subparagraph{Directional extensions}~

\begin{table}[H]
	\caption{Sokoro motion morphology}
	\label{tab:soko}
	\begin{tabular}{llll} % add l for every additional column or remove as necessary
		\lsptoprule
		forms & source gloss &  direction  & other functions  \\ %table header
		\midrule
		\textit{-ti}  &  \textsc{obl.loc}, \textsc{past}   & \textsc{vent}  &   \\
		\lspbottomrule
	\end{tabular}
\end{table}

The suffix \textit{-ti} in Sokoro can occur with a manner of motion verb in a context where it appears to function as indicating that the path of motion is toward the deictic center (ventive). Compare (\ref{ex2:soko1}) without the suffix \textit{-ti} where the motion is away from the deictic center (the speaker), and (\ref{ex2:soko2})  with the suffix where the motion is toward the deictic center. The same pattern is seen with a transitive motion verb in (\ref{ex2:soko3}) and (\ref{ex2:soko4}) .

\ea\label{ex2:soko1}
\langinfo{Sokoro verb \textit{gəɗ} ‘run’ without ventive suffix}{}{\cite[]{yarangaEsquiseGrammaticaleLanguen.d.}} \\
\gll gəɗ	talɗi	siitdu	ono \\
run	far	??	1\gl{sg} \\
\glt `Run away from me.' [French: `\textit{Cours loin de moi.}']

\ex\label{ex2:soko2}
\langinfo{Sokoro verb \textit{gəɗ} ‘run’ with ventive suffix}{}{\cite[]{yarangaEsquiseGrammaticaleLanguen.d.}} \\
\gll gəɗ-ti	ɔt-ti	galdu	(ono) \\
run-\gl{vent}	come-\gl{vent}	??	1\gl{sg}	 \\
\glt `Run toward me.' [French: `\textit{Cours vers moi.}']

\ex\label{ex2:soko3}
\langinfo{Sokoro verb \textit{əssa} ‘pull’ without ventive suffix}{}{\cite[]{yarangaEsquiseGrammaticaleLanguen.d.}} \\
\gll əssa	səəw	talɗi	galdu \\
full	wood	far	?? \\
\glt `Pull the wood away from me.' [French: `\textit{Tire le bois loin de moi.}']

\ex\label{ex2:soko4}
\langinfo{Sokoro verb \textit{əssa} ‘pull’ with ventive suffix}{}{\cite[]{yarangaEsquiseGrammaticaleLanguen.d.}} \\
\gll əssa-ti	səəw	galdu	ono \\
pull-\gl{vent}	wood	??	1\gl{sg} \\
\glt `Pull the wood toward me.' [French: `\textit{Tire le bois vers moi.}']
\z 

The ventive directional motion verb ‘come’ can occur with or without the suffix \textit{-ti} with no change in its directional meaning, as in (\ref{ex2:soko5}) and (\ref{ex2:soko6}) (Michel Karim, personal communication). Note that the use of \textit{-ti} also triggers an irregular form of the verb stem.

\ea\label{ex2:soko5}
\langinfo{Sokoro verb \textit{obiŋ} `come’ with ventive suffix}{}{\cite[]{yarangaEsquiseGrammaticaleLanguen.d.}} \\
\gll ka	óno	ot-ti	biyo	\\
\gl{rel}.\gl{m}	1\gl{sg}	come.\gl{ipfv}-\gl{vent}	yesterday	\\
\glt `Mine came yesterday.' [French: `\textit{Le mien est venu hier.}']

\ex\label{ex2:soko6}
\langinfo{Sokoro verb \textit{obiŋ} `come’ without ventive suffix}{}{\cite[]{yarangaEsquiseGrammaticaleLanguen.d.}} \\
\gll matti	ka	kaŋ	ob-e	Wəgəle \\
man		\gl{rel}.\gl{m}	\gl{dem}.\gl{m}	come.\gl{pfv}	Gogmi \\
\glt `This man came to Gogmi.' [French: `\textit{Cet homme est venu à Gogmi.}']
\z

With some non-motion verbs, the suffix appears to have a translocative function. Compare the translations of the verb \textit{seɗ-} `leave (something)' with the suffix \textit{-ti}, in (\ref{ex2:soko10}), and without it, in (\ref{ex2:soko11}). In (\ref{ex2:soko10}), the event takes place at a different location from the place of speech.

\ea\label{ex2:soko10}
\langinfo{Sokoro verb \textit{seɗ} `leave’ with ventive suffix}{}{\cite[]{yarangaEsquiseGrammaticaleLanguen.d.}} \\
\gll seɗ-da-i	pook \\
leave.\gl{imp}-3\gl{f}-\gl{vent}	chair \\
\glt `Leave the chair over there.' [French: `\textit{Laisse la chaise là-bas.}']

\ex\label{ex2:soko11}
\langinfo{Sokoro verb \textit{seɗ} `leave’ without ventive suffix}{}{\cite[]{yarangaEsquiseGrammaticaleLanguen.d.}} \\
\gll seɗ-da	pook	\\
leave.\gl{imp}-3\gl{f}	chair \\	
\glt `Leave the chair.' [French: `\textit{Laisse la chaise.}']
\z 

In other contexts, it is even less clear what the function of \textit{-ti} is, as in examples (\ref{ex2:soko7}) and (\ref{ex2:soko8}) where the suffix in the source is glossed as past tense. \citet[48]{yarangaEsquisseGrammaticaleLangue2019} describe this temporal-aspectual use of the suffix as meaning that the event occurs at some particular moment (``\textit{à un moment précis}'').

\ea\label{ex2:soko7}
\langinfo{Sokoro verb \textit{lay} `sing’ with ventive suffix}{}{\cite[]{yarangaEsquiseGrammaticaleLanguen.d.}} \\
\gll Kee ki, kee daa-g-a yeer-um kaɗi ki, mba gəəɗ kee lay-i-ti ma gaam ma? \\
2\gl{sg} \gl{top}, 2\gl{sg} kill-3\gl{f}-\gl{ipfv} mother-2\gl{sg}.\gl{m} ?? \gl{top} but song 2\gl{sg} sing-2.\gl{irr}-\gl{vent} \gl{q} none \gl{q}	\\
\glt `You there, you killed your mother. Do you want to sing or not?' \newline [French: `\textit{Toi là, tu as tué ta mère, est-ce que tu vas chanter ou non} ?']

\ex\label{ex2:soko8}
\langinfo{Sokoro verb \textit{wag} `pound’ with ventive suffix}{}{\cite[]{yarangaEsquiseGrammaticaleLanguen.d.}} \\
\gll Ka gin-ti ma? Naa wag-ti yeer	 \\
2\gl{sg} do.\gl{ipfv}-\gl{vent} \gl{q} 1\gl{sg} pound.\gl{ipfv}-\gl{vent} milet		 \\
\glt `What was I doing this morning? I was pounding millet.' \newline [French: `\textit{Qu’est-ce que tu faisais ce matin ? Je pilais le mil.}']
\z

With non-motion verbs, the suffix \textit{-ti} frequently appears in a context where translational movement is described in an immediately adjacent clause or in a preceding verb.

\subparagraph{Caused accompanied motion (CAM)}

The ventive CAM verb is \textit{aaya} `\textit{amener} [bring]' which often appears with the ventive extension: \textit{ay-ti} `bring (here)'. 

\subparagraph{Buy-sell verbs}

There are two separate verbs in Sokoro for `buy' and `sell': \textit{epir} `buy' and \textit{diging} `sell' (Michel Karim, personal communication). \citet{yarangaEsquiseGrammaticaleLanguen.d.} include two examples of the verb \textit{epir} `buy' with the translocative suffix \textit{-ti}, as in (\ref{ex2:soko14}).

\ea\label{ex2:soko14}
\langinfo{Sokoro verb \textit{epir} `buy' with ventive suffix}{}{\cite[]{yarangaEsquiseGrammaticaleLanguen.d.}} \\
\gll Kee gara epir-ti yeeri ma?	\\
2\gl{sg} want.\gl{ipfv} buy.\gl{irr}-\gl{vent} millet \gl{q}	\\
\glt `Do you want to buy some millet?' [French: `\textit{Tu veux acheter du mil} ?']
\z 

\subsection{East Chadic B4 (1/2)}

This subgroup consists of several language varieties which all respond to the exonym Barayin. The variety included here consists of the Jalkiya and Giliya lects, which are nearly identical \citep{lovestrandDialectsBarainEast2011}. In contrast, the Kodi and Jalking dialects can be considered separate languages. Only the Jalking (saka1296) variety is given a separate code in the Glottolog; neither are recognized by the Ethnologue. 

\paragraph*{Barayin (bare1279)} \label{bara}

\subparagraph{Directional verbs}~

\begin{table}[H]
	\caption{Barayin directional verbs} 
	\label{tab:baraV}
	\begin{tabular}{lll} % add l for every additional column or remove as necessary
		\lsptoprule
		DIR & verb & gloss   \\ %table header
		\midrule
		ITIVE & \textit{kóló} & `go' \\
		VENT & \textit{ājó, síi} & `come' \\
		UP & \textit{tado} & `ascend' \\
		DOWN & \textit{jàŋgó} & `descend' \\
		INTO& \textit{tōpó} & `enter' \\	
		OUT & \textit{gùsò} & `exit' \\
		\lspbottomrule
	\end{tabular}
\end{table}

\subparagraph{Directional extensions}~

There is no motion-related verbal morphology in Barayin \citep{lovestrandLinguisticStructureBarain2012}.

\subparagraph{Caused accompanied motion (CAM)}

There are several verbs in Barayin that can express directed CAM. There is a ventive-itive opposition lexicalized in the CAM verbs \textit{kóoro} `take away' and \textit{sèeró} `bring' or \textit{ajuro} `bring'. The verbs are transparently related to the motion verbs, \textit{kóló} `go' and \textit{ājó, síi} `come'. The stem-final \textit{-r} (before the infinitival suffix \textit{-o}) would appear to be a fossilized ``oblique'' suffix which had a valency-increasing function \citep{robertsObliquesMawa2016}. However, the forms of two of the directed CAM verbs show irregular phonological patterns, indicating that a lexicalization of these stems has occurred. 

More directed CAM verbs may be found in a small class of causative verbs that, in comparison with their non-causative equivalents, show an apparent stem-internal vowel alternation. This pattern is limited to a small class of verbs. Examples of this pattern include \textit{gùsò} `exit' compared with \textit{gīsō} `take out', \textit{tādó} `ascend' compared with \textit{tīdó} `take up' and \textit{jàŋgó} `descend' compared with \textit{jìŋgó} `take down'. 

\subparagraph{Buy-sell verbs}

There are two verbs in Barayin: \textit{gōrō} `buy' and \textit{bìtó} `sell, lose'. 
